\chapter{FuseSoC and the .core File}
\label{ch:fusesoc}

\chapabstract{What \FuseSoC actually does, how a \file{.core} file describes
a design, and how to read one of the lab's \file{.core} files end to end.
Ends with a checklist for writing your own.}

\section{What FuseSoC is}

\FuseSoC is a package manager and build tool for hardware sources. It solves
one problem: a design is rarely a pile of files handed straight to a
simulator. It depends on other designs (IPs), it needs a different file list
for simulation than for synthesis, and each tool (\tool{Verilator},
\tool{QuestaSim}, \tool{Design Compiler}, \dots) wants its options passed
differently. \FuseSoC lets you describe all of that once, declaratively, in a
\file{.core} file, and then resolves it: given a target (``simulate'',
``lint'', ``synthesise''), it works out which files are needed, in which
order, pulls in whatever those files themselves depend on, and calls the
right tool with the right flags.

Every design and every dependency, however small, is a \textbf{core}: a
\file{.core} file plus the sources it points at. A lab accelerator is a core.
\file{common\_cells} is a core. X-HEEP itself is a core. \FuseSoC does not
care how big a core is, only that it has a name and a \file{.core} file
describing it.

\begin{ruleofthumb}
A \file{.core} file does not contain RTL. It only says which files exist,
what they depend on, and how to run them under each target. Read it as a
table of contents plus a recipe, never as an implementation.
\end{ruleofthumb}

\section{Anatomy of a .core file}

Every \file{.core} file in these labs has the same four kinds of section, in
the same order: a header, a name and description, one or more
\kw{filesets}, and one or more \kw{targets}. Here is the whole of
\file{cordic\_accel.core} from Lab 1, the smallest real example in the
series.

\begin{corecode}[caption={cordic\_accel.core -- lab1/cordic\_accel.core (trimmed)},label={lst:fusesoc:core}]
CAPI=2:
name: isa:lab1:cordic_accel:0.1.0
description: CORDIC accelerator with an OBI data port and a reggen CSR interface

filesets:
  rtl:
    files:
      - rtl/cordic_pkg.sv
      - rtl/cordic_accel.sv
    file_type: systemVerilogSource
    depend:
      - pulp-platform.org::obi
      - pulp-platform.org::register_interface
      - pulp-platform.org::common_cells

  tb:
    files:
      - tb/tb_cordic_accel.sv
    file_type: systemVerilogSource
    depend:
      - pulp-platform.org::obi_test
      - pulp-platform.org::reg_test

targets:
  default: &default
    filesets: [rtl]

  sim:
    <<: *default
    default_tool: verilator
    filesets: [rtl, tb]
    toplevel: tb_cordic_accel
    tools:
      verilator:
        mode: binary
        verilator_options: [--timing, --trace-fst]

  lint:
    <<: *default
    default_tool: verilator
    toplevel: cordic_accel
    tools:
      verilator:
        mode: lint-only
\end{corecode}

\subsection{The header and the name}

\code{CAPI=2:} is not a comment: it is the first line \FuseSoC's parser
looks for, and it says ``this file uses Core API version 2''. Every
\file{.core} file in these labs starts with it.

The \kw{name} is a colon-separated identifier, conventionally
\code{<vendor>:<library>:<core>:<version>}: \code{isa:lab1:cordic\_accel:0.1.0}
reads as ``the ISA lab group's \code{lab1} library, core
\code{cordic\_accel}, version 0.1.0''. This is exactly what another
\file{.core} file names in its own \kw{depend} list to pull this core in;
\cref{ch:vendoring} shows the vendored IPs' names the same way
(\code{pulp-platform.org::obi} and similar).

\subsection{Filesets}

A \kw{fileset} is a named group of files that belong together for one
purpose: \code{rtl} for the synthesisable design, \code{tb} for the
testbench, and so on -- the names themselves are not special to \FuseSoC,
only convention. Each fileset lists:

\begin{itemize}
  \item \kw{files}: paths relative to the \file{.core} file itself, in
        compile order where the tool cares (packages before the modules that
        use them).
  \item \kw{file\_type}: how the tool should treat them --
        \code{systemVerilogSource}, \code{vhdlSource}, \code{cSource}, and
        so on.
  \item \kw{depend}: other cores' names. \FuseSoC resolves these recursively
        and adds their own filesets to the build, so you never list a
        dependency's files by hand.
\end{itemize}

\begin{gotcha}{A fileset is not a target}
Declaring \code{rtl:} and \code{tb:} does not build anything by itself.
Nothing runs until a \emph{target} lists which filesets it needs. It is
entirely normal for a \file{.core} file to declare a fileset that only one
target ever uses (\cref{lst:fusesoc:core}'s hypothetical synthesis-only
technology cells, for instance).
\end{gotcha}

\section{Targets: the same sources, run differently}

A \kw{target} is a named recipe: which filesets to pull in, which tool to
run, which module is the top of the design, and any tool-specific options.
The real \file{cordic\_accel.core} (\cref{ch:vendoring} onwards use the
un-trimmed file) has five: \code{default}, \code{sim} (Verilator),
\code{sim\_questa} (QuestaSim, same testbench), \code{synth} (Design
Compiler) and \code{lint} (Verilator, lint-only). They all draw from the same
two or three filesets; only the combination and the tool options change.

\begin{notebox}{The \code{\&default} / \code{<<: *default} trick}
\code{\&default} and \code{<<: *default} are plain YAML anchors and merge
keys, not \FuseSoC syntax: \code{\&default} names the \code{default:} block,
and \code{<<: *default} copies its keys into another block before
overriding some of them. It exists purely so you do not repeat
\code{filesets: [rtl]} in every target.
\end{notebox}

\kw{toplevel} names the module the tool should elaborate from; it can differ
between targets in the same file (\cref{lst:fusesoc:core}'s \code{sim}
target elaborates the testbench, \code{lint} elaborates the accelerator
itself, since there is no testbench to lint). The \kw{tools} block holds
one sub-block per tool named in \kw{default\_tool}, with whatever options
that specific tool understands -- \code{verilator\_options} for
\tool{Verilator}, \code{vlog\_options} for \tool{QuestaSim}'s compiler, and
so on. \FuseSoC does not validate these against the tool; a typo here fails
at the tool invocation, not at parse time.

\section{Running it}

\begin{shcode}
fusesoc --cores-root . --cores-root vendor run --target sim isa:lab1:cordic_accel
\end{shcode}

Reading left to right: \code{--cores-root <dir>} tells \FuseSoC where to
look for \file{.core} files, and it needs both the lab directory itself
(for \file{cordic\_accel.core}) and \file{vendor/} (for the \file{.core}
files describing the vendored IPs it depends on) -- omit either and
resolution fails with ``core not found''. \code{run --target sim} picks the
target block by name. The last argument is the core's \kw{name}. \FuseSoC
then resolves the dependency tree, writes a file list into a
\file{build/<core>/} directory, and invokes the target's \code{default\_tool}
on it.

\begin{tipbox}{Let the Makefile remember the flags}
Nobody types the \code{--cores-root} flags by hand every time; that is
exactly what the lab Makefiles are for. \Cref{ch:makefiles} shows the
wrapper.
\end{tipbox}

\section{Where the build goes}
\label{sec:fusesoc:builddir}

Every run writes into \file{build/<core-name>/<target>-<tool>/}, with colons
in the core's \kw{name} turned into underscores -- \file{isa:lab1:cordic\_accel:0.1.0}
and target \code{sim} become \file{build/isa\_lab1\_cordic\_accel\_0.1.0/sim-verilator/}.
Inside that folder: the file list \FuseSoC derived (\file{*.eda.yml}), a
generated \file{Makefile} that actually drives the tool, and every
intermediate and output file the tool produces -- for \tool{Verilator},
that includes the compiled simulation binary and, once you run it, the
waveform dump. None of it is meant to be edited or committed; delete the
whole \file{build/} tree whenever a run looks confused and re-run
\code{fusesoc} to start clean.

\begin{shcode}
fusesoc --cores-root . --cores-root vendor core-info isa:lab1:cordic_accel
fusesoc --cores-root . --cores-root vendor list-cores
\end{shcode}

\code{core-info} prints the resolved \file{.core} file -- filesets, targets
and all -- exactly as \FuseSoC parsed it, which is the fastest way to catch
a YAML typo or a missing \kw{depend} before running anything. \code{list-cores}
lists every core \FuseSoC can currently see across the given
\code{--cores-root} directories; if your own core is missing from it, the
directory holding its \file{.core} file was never passed in.

\begin{gotcha}{\code{-Wno-fatal} is not ``skip linting''}
\Cref{lst:fusesoc:core}'s real counterpart adds \code{-Wno-fatal} to the
\code{sim} target's \tool{Verilator} options, with a comment worth reading
literally: \emph{``The vendored IPs are lint-clean for their own flow, not
for ours.''} \tool{Verilator} still reports every warning; \code{-Wno-fatal}
only stops a warning from \emph{aborting} the build, so warnings in
\file{vendor/} code (built against a different lint configuration upstream)
do not block a simulation of your own module. It does not silence output
and it is not a licence to ignore warnings in the RTL you wrote yourself --
\code{lint} runs without it for exactly that reason.
\end{gotcha}

\section{Writing your own .core file}

Copying \file{cordic\_accel.core} and editing it is the right way to start;
these are the pieces to change, in order.

\begin{enumerate}
  \item \textbf{Name it.} Pick \code{<yourvendor>:<lab>:<core>:<version>}.
        This is the string every dependant, and every \code{fusesoc run},
        will use to refer to it.
  \item \textbf{List your sources in a fileset.} Compile order matters for
        SystemVerilog packages: a \code{\_pkg.sv} file goes before anything
        that imports it.
  \item \textbf{Declare dependencies by name}, not by copying their files
        into your fileset. If a dependency is not already vendored,
        \cref{ch:vendoring} is the next stop, not this file.
  \item \textbf{Add a target for the tool you actually run.} Start from the
        smallest working target (usually \code{lint}, since it needs no
        testbench) and add \code{sim} once the design elaborates cleanly.
  \item \textbf{Set \kw{toplevel}} to whatever module that target should
        elaborate -- the design for \code{lint}, the testbench for
        \code{sim}.
\end{enumerate}

\section{Checklist}
\label{sec:fusesoc:checklist}

\begin{itemize}
  \item A \file{.core} file describes files and how to run them; it never
        contains RTL itself.
  \item A fileset groups files with a type and a dependency list; a target
        picks filesets, a tool, and a toplevel.
  \item \kw{depend} names another core; \FuseSoC resolves it, you do not
        list its files.
  \item \code{fusesoc run} always needs every \code{--cores-root} that holds
        a \file{.core} file the build touches, your own directory included.
  \item When something does not build, read the target block first: wrong
        filesets or the wrong toplevel accounts for most first mistakes.
\end{itemize}
